\section{Discussion}

\subsection{A Regime-Dependent Trade-off}
The dyadic hierarchy of HEALPix has a real cost. After compression,
native UTM stays the more compressible representation at every quality
tested except the loosest, and compression
narrows the gap steadily, from $+107\%$ to $-7\%$. The operational
verdict is therefore decided by whether the residual gap exceeds the
overhead a given native layout carries, the \term{critical overhead}
$\gamma_\star$ of Table~\ref{tab:synthesis}: far above the overheads
examined here at high fidelity, $1.23$ at
$3\times10^{-3}$, where the verdict depends on the layout, and $0.93$ at
$10^{-2}$, where the \term{carrier} is smaller at equal layout overhead,
though not significantly on the ground both grids share. A
HEALPix \term{carrier} pays less for its oversampling than raw byte counts
suggest, and compression does not make it free.

\subsection{What Is a Property of the Grid}
\label{sec:grid_property}
One result is purely geometric: the raw $2.584\times$ sampling-density
ratio, which the illustrative $\gamma_{\mathrm{tile}}=1.30$ turns into a
product-architecture quantity of $1.99\times$. A second behaviour is
reproduced by geometry alone under the measured signal model: the
monotone decrease of the overhead-free \term{size ratio} with the \term{target}. An
analytic field sampled on both lattices reproduces its shape and, to
within $5$--$24\%$, its magnitude. It matches to within $7\%$ once the
native side is read through the same interpolant, and it reproduces the
crossover once the field itself passes through the \term{transfer}
(Section~\ref{sec:support_controls}). The level-19 \term{carrier} swept under
the same protocol says the same from below
(Section~\ref{sec:level_choice}). Sampling geometry is therefore
sufficient to generate this trend under the conditions tested. The
exploratory latitude dependence is consistent with HEALPix cell-shape
anisotropy, and has not been isolated as a grid property. Three
findings belong elsewhere. The crossing of parity belongs to
the \term{transfer}: it exists for a Clough--Tocher \term{carrier} and for no
other \term{transfer} tested, and it disappears when the native raster is
given that smoothing. It also belongs to the codec: with no overhead
only JPEG~XL shows it on its own, and under the illustrative overhead
four of the five fall below one and ZFP does not. A substantial part of the
transfer-and-read floor is caused by the \term{read}, since changing the
\term{read}ing operator alone removes two thirds of its mean squared error
(Section~\ref{sec:operator}). The codec ranking is a property of these
codecs on this imagery, identical on both grids, with no demonstrated
quadtree effect (Section~\ref{sec:why_quadtree}). A reader interested
in HEALPix as such should keep the first two; a product designer
should keep all five, because a product is the complete chain.
Full-sphere behaviour has not been measured: every result is on one
level-10 parent tile per scene, and a global product would need a
multi-chunk hierarchy and boundary handling across the 12 base faces of
level 0.

One architectural option follows from Section~\ref{sec:redundancy} and is
not measured here. The Clough--Tocher \term{carrier} holds no more information
than the native window it was built from, so a store that keeps the
native product and applies the \term{transfer} at read time carries no
storage penalty: the HEALPix index becomes a view over native chunks
rather than a second copy of the data. The storage arithmetic favours
this design most for the PSF-aware \term{transfer}, whose tile costs
$2.3$--$2.6$ times the \term{native square} at every \term{target}
(Section~\ref{sec:transfer_cost}). The native pixels that build one
tile are those under its footprint plus a halo of a few pixels for the
support of the kernel: a halo of four cells enlarges the footprint by
$(1032/1024)^2-1\approx1.6\%$ and four native pixels by $2.5\%$, an
order of magnitude below the illustrative $30\%$, and such a store
keeps only valid pixels. Two open questions decide whether the design
pays. A native product compressed to a given \term{target} must still
carry the small scales the deconvolution restores, without the
deconvolution amplifying the coding error more than it restores; and
the \term{transfer}, the expensive step of the chain, then runs at every
\term{read} instead of once. Both are future work. Every ratio in this
paper is about the other choice, a \term{carrier} materialised on disk,
which is what a DGGS archive buys latency and simplicity with; a hybrid
that materialises the \term{carrier} only for the levels and regions
read often sits between the two.

\subsection{What a Shared Grid Adds That This Study Does Not Count}
\label{sec:shared_grid}
Every ratio above compresses one tile, one band and one date as an
independent two-dimensional raster, the conditions of the native
product, chosen to isolate what the grid itself costs. They leave out
the reason a DGGS is used. On HEALPix every product, whatever its
sensor, band or date, lands on the same cells with the same
identifiers, and the nested index makes any block of cells one
contiguous slice. A codec can then exploit the correlation between
dates, bands and sensors, where native products would first have to be
reprojected onto one another. For a single tile and sensor the native
stack offers the same alignment in time; the difference appears across
sensors, across UTM zones and in the overlaps of neighbouring tiles. The
gain measured here, a \term{carrier} that is competitive only at
aggressive compression, prices the geometry alone and none of that
coherence, which is the next measurement. Gridded model output is its
natural first case. It is already a cell average, so there is no
instrument response to model, though a model that runs on another grid
must still be remapped, with an error that conserving the integral
does not remove; its dominant redundancy is temporal; and it is global,
where the equal-area cells and the locality of the nested index matter
most.

\subsection{Reproducibility}
\label{sec:reproducibility}
The figures quoted in the text are produced by scripts from the stored
results, and an assertion harness recomputes them and fails when the
text does not carry a value it computes. The
sizes reported are those of files written, closed and re-read by an
independent decoder (Appendix~\ref{app:codec_audit}). Each sweep
records the codec library versions and the tagged checkout of the
shared resampling package~\cite{delouis2026resample} it ran against,
and the harness refuses a run whose instrument response falls outside
the band of the sensor. The code is at \repourl; that tag is archived
at \codedoi. The native windows, the stored tiles, every result file
and the per-region values this paper does not print are deposited at
\datadoi. The large-scale trial of Appendix~\ref{app:large_scale_trial}
is the one exception: it ran a conversion chain on an object store,
and its scripts and notebook are in the repository without their
inputs.

\subsection{Current Limitations}
\label{sec:limitations}
The main limits are stated with each result. The crossover holds for a
\term{carrier} filled by Clough--Tocher and scored against its own
smooth tile. The PSF-aware \term{transfer} is priced at
$\lambda=10^{-3}$, not at the $\lambda=0.1$ of the conversion chain of
Appendix~\ref{app:large_scale_trial}. The evidence is Sentinel-2 B04 at
\SI{10}{m}, one acquisition and one tile per region, and the overhead of
the native product is a parameter, not a measurement.
Appendix~\ref{app:limitations} lists every limit with its section.
